\subsection{扩张模与挠模的有限性}

设 M 是左 A-模，I 是滤过预序集， \((\mathbf{N}_{i}, u_{ji})\) 是关于 I 的左 A-模归纳系， \(N = \varinjlim N_{i}\) 为其归纳极限， \(u_{i}: N_{i} \to N, i \in I\) 为典范映射。则 \((\operatorname{Ext}_{\mathbf{A}}(\mathbf{M}, \mathbf{N}_{i}), \operatorname{Ext}_{\mathbf{A}}(1_{\mathbf{M}}, u_{ji}))\) 是 k-模的归纳系， \((\operatorname{Ext}_{\mathbf{A}}(1_{\mathbf{M}}, u_{i}))\) 是映射的归纳系，其归纳极限是分次 k-模的同态，称为典范同态

\[
\varinjlim_ {i \in I} \operatorname{Ext} _ {\mathbf {A}} (\mathbf {M}, \mathbf {N} _ {i}) \rightarrow \operatorname{Ext} _ {\mathbf {A}} \left(\mathbf {M}, \varinjlim_ {i \in I} \mathbf {N} _ {i}\right).\tag{8}
\]

命题 5. --- 若 A 是左诺特的且 M 是有限生成 A-模，则典范同态 (8) 是双射的。

事实上，设（X，第 53 页，命题 6）\(p: L \to M\) 为 M 的自由分解，使得 \(L_{n}\) 对每个 \(n\) 都是有限生成的。从 \(k\) -复形

\[
u: \varinjlim \operatorname{Homgr} _ {\mathbf {A}} (\mathrm{L}, \mathrm{N} _ {i}) \rightarrow \operatorname{Homgr} _ {\mathbf {A}} (\mathrm{L}, \varinjlim \mathrm{N} _ {i})
\]

是双射的，因此同态

\[
\varinjlim \mathrm{H} (\operatorname{Homgr} _ {\mathrm{A}} (\mathrm{L}, \mathrm{N} _ {i})) \rightarrow \mathrm{H} (\operatorname{Homgr} _ {\mathrm{A}} (\mathrm{L}, \varinjlim \mathrm{N} _ {i}))
\]

推出同态 \(\mathrm{H}(\mathrm{Homgr}(1,u_{i}))\) （X，第 28 页，命题 1）。然后由命题 2（X，第 103 页）和定理 1（X，第 100 页）得出结论。

命题 6. --- 设 B 是环，N 是 (A, B)-双模，且为诺特 B-模（相应地，有限长度）。

\begin{enumerate}
\def\labelenumi{\alph{enumi})}
\item
  假设 A 是右诺特的，设 M 是有限生成右 A-模。则 B-模（X，第 81 页）Tor \(_{n}^{A}\) (M, N) 是诺特的（相应地，有限长度的）。
\item
  假设 A 是左诺特的，设 M 是有限生成左 A-模。则 B-模（X，第 98 页）Ext \(_{A}^{n}\) (M, N) 是诺特的（相应地，有限长度的）。
\end{enumerate}

选取自由分解 \(p: \mathbf{L} \to \mathbf{M}\) ，使得每个 A-模 \(\mathbf{L}_n\) 都是有限生成的（X，第 53 页，命题 6），并设 C 为 B-模复形 \(\mathbf{L} \otimes_{\mathbf{A}} \mathbf{N}\) （情形 \(a\) )）或 \(\operatorname{Homgr}_{\mathbf{A}}(\mathbf{L}, \mathbf{N})\) （情形 \(b\) )）。每个 B-模 \(\mathbf{C}_n\) 同构于 N 的有限个拷贝的乘积，因此是诺特的（相应地，有限长度的）；从而模 \(\mathrm{H}_n(\mathbf{C})\) 亦然。现在，根据 X，第 100 页，定理 1，这些模同构于 \(\operatorname{Tor}_n^{\mathbf{A}}(\mathbf{M}, \mathbf{N})\) （情形 \(a\) )）或 \(\operatorname{Ext}_{\mathbf{A}}^{-n}(\mathbf{M}, \mathbf{N})\) （情形 \(b\) )）。

推论. --- 设 \(\rho: A \to B\) 为诺特交换环的同态，M 为有限生成 A-模，N 为 B-模。若 N 是有限生成（相应地，有限长度）B-模，则 B-模 \(\operatorname{Tor}_{n}^{A}(M, N)\) 与 \(\operatorname{Ext}_{A}^{n}(M, N)\) 也是。

